\documentclass[12pt]{article}
\usepackage[utf8]{inputenc}
\usepackage{amsmath,amssymb,amsthm}
\usepackage{geometry}
\geometry{margin=1in}

\newtheorem{theorem}{Theorem}
\newtheorem{lemma}[theorem]{Lemma}
\newtheorem{corollary}[theorem]{Corollary}

\title{Problem 113: Non-bouncy Numbers}
\author{Project Euler Solution}
\date{}

\begin{document}
\maketitle

\section{Problem Statement}

A positive integer is \textbf{increasing} if its digits are non-decreasing left to right, \textbf{decreasing} if non-increasing, and \textbf{bouncy} otherwise. How many numbers below $10^{100}$ are not bouncy?

\section{Mathematical Development}

\begin{theorem}[Increasing count]
The number of positive increasing integers with at most $n$ digits is
\[
I(n) = \binom{n+9}{9} - 1.
\]
\end{theorem}

\begin{proof}
A non-decreasing sequence of $n$ digits from $\{0,\ldots,9\}$ (with leading zeros) corresponds bijectively to a multiset of size~$n$ from a 10-element alphabet. By stars and bars, the count is $\binom{n+9}{9}$. Subtracting the all-zeros sequence (representing 0) gives the result.
\end{proof}

\begin{theorem}[Decreasing count]
The number of positive decreasing integers with at most $n$ digits is
\[
D(n) = \binom{n+10}{10} - 1 - n.
\]
\end{theorem}

\begin{proof}
For each length $k$, non-increasing sequences from $\{0,\ldots,9\}$ number $\binom{k+9}{9}$, of which one is all zeros. Thus the count of positive decreasing numbers of exactly $k$ digits is $\binom{k+9}{9} - 1$. Summing over $k = 1,\ldots,n$ and applying the hockey-stick identity $\sum_{k=1}^{n}\binom{k+9}{9} = \binom{n+10}{10} - 1$ yields
\[
D(n) = \binom{n+10}{10} - 1 - n.
\]
\end{proof}

\begin{lemma}[Overlap]
The number of positive integers with at most $n$ digits that are both increasing and decreasing is $9n$.
\end{lemma}

\begin{proof}
Such numbers are repdigits. For each length $k \in \{1,\ldots,n\}$ and digit $d \in \{1,\ldots,9\}$, there is exactly one repdigit, giving $9n$ in total.
\end{proof}

\begin{theorem}[Non-bouncy formula]
The number of non-bouncy positive integers below $10^n$ is
\[
\operatorname{NonBouncy}(n) = \binom{n+9}{9} + \binom{n+10}{10} - 10n - 2.
\]
\end{theorem}

\begin{proof}
By inclusion--exclusion:
\begin{align*}
\operatorname{NonBouncy}(n) &= I(n) + D(n) - 9n \\
&= \bigl[\tbinom{n+9}{9} - 1\bigr] + \bigl[\tbinom{n+10}{10} - 1 - n\bigr] - 9n \\
&= \binom{n+9}{9} + \binom{n+10}{10} - 10n - 2. \qedhere
\end{align*}
\end{proof}

\begin{corollary}[Asymptotic vanishing]
$\operatorname{NonBouncy}(n)/(10^n - 1) \to 0$ as $n \to \infty$, since $\operatorname{NonBouncy}(n) = O(n^{10})$.
\end{corollary}

\section{Algorithm}

Evaluate $\binom{109}{9} + \binom{110}{10} - 1002$ using big-integer arithmetic (two binomial coefficients, each requiring $O(10)$ multiplications).

\section{Complexity Analysis}

\begin{itemize}
\item \textbf{Time:} $O(1)$ arithmetic operations (two binomial coefficients with constant-size lower parameters).
\item \textbf{Space:} $O(1)$ big integers.
\end{itemize}

\section{Answer}

\[
\boxed{51161058134250}
\]

\end{document}
